% SOURCE_AUTHORITY: sga5_fr_workpass.tex, lines 8649--8663 (LF-normalized unit).
% SOURCE_SHA256: 791F4EFFC5E02832D5D77ED03518C8156D6F07E4C8238B03545DB93D883FBB28
\subsubsection*{1.4.5}
Seja \(F\) um \(\mathbb Z_\ell\)-feixe construtível, e denotemos por \(p\) o funtor canônico
\[
\mathbb Z_\ell\text{-}\fc(X)\longrightarrow \mathbb Q_\ell\text{-}\fc(X).
\]
O funtor composto de \(p\) e o funtor \(G\mapsto F\otimes G\) torna inversíveis as multiplicações por elementos de \(\mathbb Z_\ell\) e define, portanto, um funtor
\[
\mathbb Q_\ell\text{-}\fc(X)\longrightarrow \mathbb Q_\ell\text{-}\fc(X),
\]
que é exato à direita. Aplicando o mesmo raciocínio ao primeiro argumento, define-se assim um bifuntor exato à direita, chamado produto tensorial:
\[
\mathbb Q_\ell\text{-}\fc(X)\times \mathbb Q_\ell\text{-}\fc(X)
\longrightarrow \mathbb Q_\ell\text{-}\fc(X).
\]
Se \(X\) é conexo e \(F\) e \(G\) são dois \(\mathbb Q_\ell\)-feixes construtíveis constantes torcidos que correspondem a representações de \(\pi_1(X)\) sobre os espaços vetoriais \(V\) e \(W\), respectivamente, o produto tensorial \(F\otimes G\) permanece localmente constante e corresponde à representação diagonal de \(\pi_1(X)\) sobre \(V\otimes W\).
